\documentclass[11pt]{article}
\usepackage[a4paper,margin=1.9cm]{geometry}
\usepackage[T1]{fontenc}
\usepackage{textcomp}
\usepackage[spanish]{babel}
\usepackage{amsmath,amssymb}
\usepackage{enumitem}
\usepackage{tikz}
\usetikzlibrary{calc,arrows.meta,patterns,decorations.pathmorphing}
\usepackage{xcolor}
\usepackage{multicol}
\usepackage{parskip}
\usepackage{lmodern}
\renewcommand{\familydefault}{\sfdefault}

\definecolor{unalblue}{RGB}{31,73,124}
\definecolor{unalblue2}{RGB}{70,120,180}

\newlist{opciones}{enumerate}{1}
\setlist[opciones]{label=\textbf{(\alph*)},itemsep=1pt,topsep=2pt,leftmargin=1.8em,font=\normalfont}
\setlist[enumerate,1]{label=\textbf{\arabic*.},itemsep=6pt,topsep=4pt,leftmargin=1.6em,font=\normalfont}
\setlist[enumerate,2]{label=\textbf{\alph*)},itemsep=4pt,topsep=3pt,leftmargin=1.6em,font=\normalfont}
\setlist[enumerate,3]{label=\textbf{\roman*)},itemsep=2pt,topsep=2pt,leftmargin=1.4em,font=\normalfont}

\newcommand{\vect}[2]{\begin{pmatrix}#1\\#2\end{pmatrix}}
\newcommand{\vectiii}[3]{\begin{pmatrix}#1\\#2\\#3\end{pmatrix}}

\newcommand{\examheader}{%
\noindent\begin{tikzpicture}
\node[fill=unalblue, text width=\dimexpr\linewidth-16pt\relax, inner sep=8pt, rounded corners=3pt, align=left, text=white] (hdr) {%
  \begin{minipage}[c]{0.66\linewidth}
    {\Large\bfseries \examsubject}\\[2pt]
    {\small \examtype\ \textbullet\ \textcolor{white!85!unalblue}{\examsubtitle}}
  \end{minipage}\hfill
  \begin{minipage}[c]{0.28\linewidth}
    \raggedleft{\scriptsize Escuela de Matem\'aticas\\[-1pt] Universidad Nacional\\[-1pt] de Colombia, Sede Medell\'in}
  \end{minipage}%
};
\end{tikzpicture}\\[7pt]
}
\newcommand{\examfooter}{%
  \vfill
  \rule{\linewidth}{0.4pt}\\[1pt]
  {\scriptsize\textit{Transcripci\'on digital --- MathUNAL}\hfill\textcolor{unalblue}{\textbf{\thepage}}}
}
\pagestyle{empty}

\newcommand{\examsubject}{CÁLCULO EN VARIAS VARIABLES}
\newcommand{\examtype}{Segundo Parcial}
\newcommand{\examsubtitle}{Semestre 02-2013, 2 de Noviembre de 2013}

\begin{document}

\examheader

\vspace{4pt}
\noindent
\begin{minipage}[t]{0.55\linewidth}
Duraci\'on: 1 hora y 40 minutos
\end{minipage}%
\hfill
\begin{minipage}[t]{0.4\linewidth}
Nombre: \dotfill\\[6pt]
Doc. Identidad: \dotfill\\[6pt]
Profesor: \dotfill \hfill Grupo: \dotfill
\end{minipage}

\vspace{6pt}
\noindent\textbf{Instrucciones.} El examen tiene una duraci\'on de 1 hora y 40 minutos. \textbf{No se admiten preguntas durante el examen.} No est\'a permitido el uso de calculadora, ni tel\'efonos m\'oviles, ni cualquier otro tipo de aparato electr\'onico. Todos los puntos de procedimiento deben estar debidamente justificados.

\begin{enumerate}
\item \textbf{(20\%)} \textbf{Selecci\'on m\'ultiple.}
\begin{enumerate}[label=\alph*)]
\item \textbf{(10\%)} Considere los siguientes dos enunciados:

\textbf{(I)} Sea $f:\mathbb{R}^2\to\mathbb{R}$ una funci\'on de clase $C^1$. Si $(2,1)$ es un punto cr\'itico de $f$ y $f_{xx}(2,1)f_{yy}(2,1) < [f_{xy}(2,1)]^2$, entonces $f$ tiene un punto silla en $(2,1)$.

\textbf{(II)} La integral $\displaystyle\int_0^{2\pi}\int_0^2\int_r^2 1\,dz\,dr\,d\theta$ representa el volumen encerrado por el cono $z=\sqrt{x^2+y^2}$ y el plano $z=2$.

Podemos afirmar que:
\begin{opciones}
\item Los enunciados (I) y (II) son falsos.
\item Los enunciados (I) y (II) son verdaderos.
\item El enunciado (I) es falso y el enunciado (II) es verdadero.
\item El enunciado (I) es verdadero y el enunciado (II) es falso.
\end{opciones}

\item \textbf{(10\%)} El valor de la integral $\displaystyle\iiint_B e^{x+y+z}\,dV$ sobre el cubo $B=[0,1]\times[0,1]\times[0,1]$ es
\begin{opciones}
\item $(e^2-1)^3$
\item $e^3$
\item $(e+1)^3$
\item $(e-1)^3$
\end{opciones}
\end{enumerate}

\vspace{10pt}
\item \textbf{(35\%)}
\begin{enumerate}[label=\alph*)]
\item \textbf{(18\%)} Use multiplicadores de Lagrange para determinar cu\'al es el \'area m\'axima que puede tener un rect\'angulo si la longitud de su diagonal es 2.
\item \textbf{(17\%)} Eval\'ue la integral $\displaystyle\int_0^1\int_0^{\cos^{-1}y} e^{\operatorname{sen}x}\,dx\,dy$. Dibuje y describa la regi\'on de integraci\'on.
\end{enumerate}
\end{enumerate}

\examfooter
\newpage

\examheader

\begin{enumerate}
\setcounter{enumi}{2}
\item \textbf{(25\%)} Sea $E$ el s\'olido que se encuentra por dentro de la superficie $x^2+y^2+z^2=2z$ y debajo de la superficie $z=\sqrt{3(x^2+y^2)}$. Si en cada punto $(x,y,z)\in E$ la densidad es igual a la distancia del punto al origen, halle la masa de $E$. Dibuje y describa la regi\'on de integraci\'on.

\vspace{10pt}
\item \textbf{(20\%)} Use el teorema de cambio de variables para evaluar la integral $\displaystyle\iint_D \operatorname{sen}(9x^2+4y^2)\,dA$, donde $D$ es la regi\'on del primer cuadrante acotada por la elipse $9x^2+4y^2=1$. Dibuje y describa las regiones de integraci\'on.
\end{enumerate}

\examfooter

\end{document}
